\setrunninghead{引子\q “不上班，也有钱”，真的能做到！}
\section*{引子\\ “不上班，也有钱”，真的能做到！}
\addcontentsline{toc}{section}{\protect\textbf{引子\q “不上班，也有钱”，真的能做到！}}

{\kaishu 钉大（银行螺丝钉的粉丝对其的昵称），你好，我今年25岁，有可能在45岁左右实现财务自由吗？}

经常有朋友问螺丝钉（银行螺丝钉的自称）这个问题，这也是很多朋友的心声，我们都想实现财务自由。

但怎样才能实现财务自由呢？需要攒够多少钱呢？

\subsubsection{40岁前就退休，和家人一起环游世界}

首先，螺丝钉给大家分享一个真实的故事。

{\kaishu 在中国宝岛台湾，有这样一个姑娘，她在10年前第一次跟男朋友去泰国“穷游”的时候，就定下了一个目标：要在40岁之前退休，和男朋友一起环游世界。}

{\kaishu 那么，怎么样才能实现这个目标呢？}

{\kaishu 这个姑娘设想了很多种方案，一开始她想通过“边旅游边打工”的方式实现，她老公也就是曾经的男朋友，甚至想过从事潜水教练、英语老师等职业。}

{\kaishu 可是在后来查找资料的过程中，这对夫妻认识到，可以通过投资实现40岁前退休的目标。于是他们确定了一个路线：在10年内，攒下足够的资产，实现财务自由。}

{\kaishu 后来，他们真的做到了。}

他们是怎么实现的呢？主要有三点：节流、开源、依靠基金分红建立被动收入来源。

先说节流。这对夫妻为了攒钱极为节省，蜜月旅行选择了徒步爬山，只花了几百元，平日从不乱买零食饮料，总是自己种菜做饭。他们在收入提高后，还制定了一个规则：每年只花掉收入的10\%，将剩下的90\%全部存下来，定期投到指数基金上。

再说开源。这对夫妻刚开始上班的时候，家庭月收入在15000元人民币左右。但是通过努力工作，10年后当他们准备退休的时候，月收入已经是刚开始工作时的3～4倍了，是5万～6万元人民币（包括工资收入和副业收入）。

最后说投资。一开始的时候他们一年投入大约8万元到指数基金，到了后期他们一年投入到指数基金的资金大概有五六十万元。

一晃10年就过去了，到了2013年，这对夫妻用近乎“自虐”的方式攒下了大约价值650万元人民币的指数基金，每年依靠指数基金的分红，就轻松实现了财务自由，达成了40岁前退休的目标，开始环游世界。

\subsubsection{从这对夫妻身上，我们能学到什么}

可能很多人会觉得这对小夫妻“太过分、太小气”，但他们不是“葛朗台”式的抠门。

葛朗台是为了抠而抠，而这对小夫妻是有自己的理想。他们愿意为了自己的理想忍受10年的辛苦生活，之后就可以享受自己想要的生活了。这样有理想和行动力的人是很可敬的。

虽然这样自虐的省钱方式，螺丝钉并不鼓励大家去模仿。但从这对小夫妻身上，我们可以借鉴以下经验。

\begin{itemize}
\item 认真工作，用双手创造价值，把自己打造成“获取稳定提升现金流”的资产，同时用保险保驾护航，这是我们的防御武器。
\item 将现金流定投到低估值的指数基金上，依靠低估值的指数基金来放大收益，这是我们的进攻武器。
\end{itemize}

这个双核制的定投体系，是最适合大多数人的。

认真工作，需要大家自己努力。而基金定投的技巧，正是本书想要介绍给大家的。

\subsubsection{实现财务自由最靠谱的方法}

说到这里，螺丝钉再给大家讲个故事。

{\kaishu 有一年春节放假的时候，螺丝钉跟两个朋友约吃饭。结果一个朋友出国旅游，另一个要加班，都没有约成。}

{\kaishu 加班的那个朋友A，工作很努力，表现也很好，单位还不错，加班给三薪，所以A也没有什么怨言。多加几天班多拿点钱，还可以接受。}

{\kaishu 另一个朋友B，原来在国企上班，后来辞职自己开了个小工作室，做设计。前几年他用赚的钱买了几套商铺，所以现在的收入是“工作收入+商铺租金”，收入还不错，经常带着家人出去旅游。他还打算学习投资，投资一些股票类品种。}

他们两个人赚的钱目前差不多，但如果不考虑对个人价值的追求，很明显朋友B过得更潇洒自在一些。

这是为什么呢？

朋友A的收入主要来自自己的人力资产：通过把自己的劳动力出租给公司，公司支付工资收入。出租自己人力资产的这段时间用来给公司创造价值，自然就干不了其他的事情。

朋友B虽然也出租自己的人力资产，但他会把所得逐渐变成资产——能收取租金的商铺。出租商铺也只需要定期收租就好，所以相对自由很多，越做越轻松。

当然能产生现金流的资产有很多，不仅仅是商铺，比如像那对中国台湾小夫妻一样持有大量的指数基金也可以收到很多现金流，道理是一样的。

出租的商铺或者指数基金等都是能产生现金流的资产，而且这种现金流的产生不需要占用你的时间。当这些资产产生的现金流能满足你所需的时候，你就不需要再出租自己的人力资产，腾出来的时间可以享受生活，或者做其他有价值的事情。这就是财务自由了。

想要变得富有，我们得攒资产；想要变得自由，我们也得攒资产。

可能大多数人未必能像那对中国台湾夫妻一样可以在10年内实现财务自由，但是通过提高自己的收入水平，学习定投的理论和方法，不断增加投资指数基金的资产比例，可以大大缩短实现财务自由的时间。原本需要工作三四十年，现在工作20年甚至只工作10年就可以实现财务自由了。

努力工作，用自己的劳动力去换资产，做一个小“资本家”，这是大多数人实现财务自由最好的方式。

我相信，我们终将富有！
